%% ===================================================================================
%% Capítulo 8 -- Resultados e Discussão
%% Parte 1 (IC) é factual, extraída do relatório de IC. As demais partes são TABELAS E
%% ROTEIROS a preencher pelo Arquiteto após a execução dos testes do Capítulo 7.
%% Regra: nunca inventar resultados. Célula sem medição permanece com "---".
%% ===================================================================================
\chapter{Resultados e Discussão}\label{cap:resultados}

Este capítulo apresenta os resultados obtidos na Iniciação Científica (Seção~\ref{sec:res-ic}), os resultados da consolidação do monólito (Seções~\ref{sec:res-requisitos} e~\ref{sec:res-metricas}), o atendimento aos objetivos (Seção~\ref{sec:res-objetivos}) e a discussão, incluindo as ameaças à validade (Seções~\ref{sec:res-discussao} e~\ref{sec:res-ameacas}).

\section{Resultados da Iniciação Científica}\label{sec:res-ic}

O relatório de IC \cite{reis2026ic} registra que os objetivos propostos foram atingidos por meio de uma abordagem de desenvolvimento incremental, e destaca as seguintes conquistas técnicas:

\begin{itemize}
  \item a implementação de um editor Markdown capaz de renderizar com precisão elementos essenciais da escrita científica, como fórmulas em \LaTeX{} e blocos de código com realce de sintaxe;
  \item o desenvolvimento de uma interface de personalização visual dinâmica, que permite gerenciar a estética do site por variáveis CSS persistidas em banco de dados, com reflexo instantâneo por ISR;
  \item a garantia da perenidade do conteúdo digital, uma vez que a arquitetura baseada em arquivos portáteis assegura que o pesquisador detenha a posse integral dos seus dados.
\end{itemize}

O mesmo relatório apresenta como limitações o prazo de execução --- que restringiu a modelagem da interface a protótipos de baixa fidelidade --- e o sistema de autenticação, que, embora seguro e funcional para a gestão básica, ainda poderia ser expandido para múltiplos níveis de permissão em cenários colaborativos.

\section{Cobertura dos Requisitos}\label{sec:res-requisitos}

A Tabela~\ref{tab:res-rf} registra a situação final de cada requisito funcional e a evidência que a sustenta (caso de teste aprovado, captura de tela ou trecho de código). A situação assume um dos valores: \emph{atendido}, \emph{parcial}, \emph{não atendido} ou \emph{adiado}. Até a conclusão da implementação, todas as células permanecem sem preenchimento.

\begin{longtable}{@{}L{1.4cm}L{4.4cm}L{2.4cm}L{5.4cm}@{}}
\caption{Situação dos requisitos funcionais}\label{tab:res-rf}\\
\toprule
\textbf{ID} & \textbf{Requisito} & \textbf{Situação} & \textbf{Evidência}\\
\midrule
\endfirsthead
\toprule
\textbf{ID} & \textbf{Requisito} & \textbf{Situação} & \textbf{Evidência}\\
\midrule
\endhead
\bottomrule
\endfoot
RF01 & Autenticação & --- & ---\\
RF02 & Estrutura hierárquica & --- & ---\\
RF03 & Ordenação & --- & ---\\
RF04 & Personalização & --- & ---\\
RF05 & Criar documento & --- & ---\\
RF06 & Criar coleção & --- & ---\\
RF07 & Editar e excluir & --- & ---\\
RF08 & Editor Markdown & --- & ---\\
RF09 & Fórmulas matemáticas & --- & ---\\
RF10 & Realce de sintaxe & --- & ---\\
RF11 & Markdown estendido & --- & ---\\
RF12 & \en{Sidebar} hierárquica & --- & ---\\
RF13 & Barra superior & --- & ---\\
RF14 & Página de documento & --- & ---\\
RF15 & Rascunho local & --- & ---\\
RF16 & Publicar versão & --- & ---\\
RF17 & Invalidação de cache & --- & ---\\
RF18 & Identidade do site & --- & ---\\
RF19 & Paleta de cores & --- & ---\\
RF20 & Encerrar sessão & --- & ---\\
RF21 & Autorização de escrita & --- & ---\\
RF22 & \en{Bootstrap} do administrador & --- & ---\\
RF23 & Configuração inicial & --- & ---\\
RF24 & Rodapé configurável & --- & ---\\
RF25 & Metadados de SEO & --- & ---\\
RF26 & Estado de publicação & --- & ---\\
\end{longtable}
\fonte{Elaborado pelos autores.}

\vitrinependencia{Arquiteto}{Preencher a Tabela de cobertura de requisitos ao final da fase F6, usando somente evidências existentes (casos de teste aprovados, capturas, trechos de código). Nenhum valor deve ser estimado.}

\section{Métricas de Qualidade}\label{sec:res-metricas}

A Tabela~\ref{tab:res-metricas} confronta as metas dos requisitos não funcionais com as medições obtidas no Capítulo~\ref{cap:testes}.

\begin{table}[htb]
\caption{Metas e medições dos requisitos não funcionais}\label{tab:res-metricas}
\small
\begin{tabularx}{\linewidth}{@{}L{1.4cm}YL{2.5cm}L{2.5cm}L{1.9cm}@{}}
\toprule
\textbf{ID} & \textbf{Métrica} & \textbf{Meta} & \textbf{Medido} & \textbf{Atende?}\\
\midrule
RNF02 & LCP (percentil 75, perfil móvel simulado) & $\leq$ 2,5 s & --- & ---\\
RNF02 & Pontuação de desempenho no Lighthouse & $\geq$ 90 & --- & ---\\
RNF05 & Pontuação de acessibilidade & $\geq$ 90 e WCAG AA & --- & ---\\
RNF04 & Pontuação SUS média & $\geq$ 68 & --- & ---\\
RNF04 & Taxa de conclusão das tarefas T1 a T6 & $\geq$ 80\% & --- & ---\\
RNF03 & Vulnerabilidades críticas em dependências & 0 & --- & ---\\
RNF01 & Comandos da instalação limpa (Apêndice~\ref{ap:implantacao}) & Informativo & --- & ---\\
RNF08 & Erros de \en{lint} / Falhas de \en{build} & 0 / 0 & --- & ---\\
--- & Casos de teste aprovados / total & Essenciais: 100\% & --- & ---\\
\bottomrule
\end{tabularx}
\fonte{Elaborado pelos autores. Medições a preencher ao término da verificação.}
\end{table}

\section{Atendimento aos Objetivos}\label{sec:res-objetivos}

\begin{table}[htb]
\caption{Atendimento aos objetivos específicos}\label{tab:res-objetivos}
\small
\begin{tabularx}{\linewidth}{@{}L{1.3cm}YL{3.6cm}L{2.0cm}@{}}
\toprule
\textbf{Objetivo} & \textbf{Resumo} & \textbf{Evidência} & \textbf{Situação}\\
\midrule
OE1 & Publicação independente por clonagem e implantação simples (imagem Docker em VPS) & Apêndice~\ref{ap:implantacao}; CT29 & ---\\
OE2 & Editor Markdown com fórmulas e realce de sintaxe & RF08 a RF11; CT15 a CT19 & ---\\
OE3 & Personalização visual sem novo \en{build} & RF19; CT23, CT24 & ---\\
OE4 & Arquitetura de dados portátil (Prisma e SQLite) & Seção~\ref{sec:arq-dados}; CT28, CT30 & ---\\
OE5 & Autenticação segura (better-auth) & RF01, RF20 a RF22; CT01 a CT06 & ---\\
OE6 & Monólito modular documentado e rastreável & Capítulos~\ref{cap:requisitos} e~\ref{cap:arquitetura}; Apêndice~\ref{ap:adr} & ---\\
OE7 & Processo DDD com agentes e rastreabilidade & Seção~\ref{sec:processo-ddd}; histórico do Git & ---\\
OE8 & Verificação e validação do sistema & Capítulo~\ref{cap:testes} & ---\\
\bottomrule
\end{tabularx}
\fonte{Elaborado pelos autores.}
\end{table}

\section{Discussão}\label{sec:res-discussao}

\subsection{Comparação com os Trabalhos Relacionados}

A discussão deve retomar a Tabela~\ref{tab:trabalhos-relacionados} e examinar, com base nas medições, em que critérios (C1 a C5) o Vitrine de fato se diferencia e em quais se aproxima ou fica aquém das categorias comparadas.

\subsection{O Processo de Desenvolvimento com Agentes}

Esta subseção avalia o método \en{Documentation-Driven Development} com agentes. Devem ser respondidas, com base no histórico do Git e no mural: quantos pedidos de integração foram aprovados de primeira e quantos foram vetados, e por quais motivos; quantas vezes o código divergiu da documentação e como a divergência foi detectada; quantas atualizações de documentação foram solicitadas pelos desenvolvedores; e qual o esforço de revisão do Arquiteto.

\vitrinependencia{Arquiteto}{Coletar, do histórico do Git e do mural, os indicadores do processo (pedidos aprovados e vetados, motivos de veto, divergências detectadas) e analisá-los nesta subseção; os indicadores só podem ser calculados após as primeiras integrações.}

\section{Ameaças à Validade}\label{sec:res-ameacas}

\begin{itemize}
  \item \textbf{Validade de construção:} as métricas de desempenho e de usabilidade são aproximações das qualidades desejadas; o LCP medido em ambiente simulado pode diferir do experimentado por usuários reais;
  \item \textbf{Validade interna:} o autor do sistema também conduz a avaliação, o que pode enviesar a interpretação; mitiga-se com casos de teste automatizados, critérios objetivos e evidências registradas;
  \item \textbf{Validade externa:} o estudo de usabilidade envolve poucos participantes e um contexto institucional específico; os resultados não devem ser generalizados para outros perfis ou áreas;
  \item \textbf{Validade de conclusão:} o volume de dados é pequeno para testes estatísticos; os resultados são descritos de modo qualitativo e com estatística descritiva;
  \item \textbf{Variabilidade dos agentes:} o comportamento de agentes de IA varia entre execuções e versões de modelos; o processo depende de revisão humana e de critérios de aceite objetivos, e os resultados do processo podem não se repetir integralmente.
\end{itemize}

\section{Limitações}

São limitações conhecidas desta versão: administrador único, sem papéis (\decisao{D-13}); ausência de histórico de versões do conteúdo, que se limita ao estado de rascunho ou publicado (\decisao{D-15}); ausência de busca textual e de envio de mídia (\decisao{D-10}); implantação restrita a ambientes com disco persistente, como uma VPS com Docker, sem suporte a plataformas \emph{serverless} (\decisao{D-04}); e o alcance da validação com usuários.

\section{Síntese do Capítulo}

O capítulo consolida os resultados factuais da IC e estrutura a apresentação dos resultados do monólito, com tabelas de cobertura de requisitos, de metas de qualidade e de atendimento aos objetivos. Os valores serão preenchidos após a verificação, sem estimativas.
